\section*{अध्याय \ref{ch:g2:water-on-earth} --- पृथ्वी पर पानी}

\begin{solution}{exo:g2:water-on-earth:1}
\omterm{def:g2:water-on-earth:salt-fresh}{खारा पानी}: \omterm{def:g2:water-on-earth:ocean}{समुद्र}, \omterm{def:g2:water-on-earth:ocean}{महासागर}। \omterm{def:g2:water-on-earth:salt-fresh}{मीठा पानी}: \omterm{def:g1:landscapes-around-us:water}{नदी}, बारिश, \omterm{def:g1:landscapes-around-us:landform}{पर्वतों} की \omterm{def:g1:landscapes-around-us:water}{झील}, \omterm{def:g1:landscapes-around-us:landform}{पर्वत}
पर की बर्फ़।
\end{solution}

\begin{solution}{exo:g2:water-on-earth:2}
क्योंकि उसकी सतह का ज़्यादातर भाग \omterm{def:g2:water-on-earth:ocean}{महासागरों} के नीले पानी से ढका है:
अंतरिक्ष से देखने पर वह नीली दिखती है।
\end{solution}

\begin{solution}{exo:g2:water-on-earth:3}
$100 - 71 = 29$ वर्ग ज़मीन हैं।
\end{solution}

\begin{solution}{exo:g2:water-on-earth:4}
\omterm{def:g2:water-on-earth:ocean}{महासागर} बड़ा है: \omterm{def:g2:water-on-earth:ocean}{समुद्र} \omterm{def:g2:water-on-earth:salt-fresh}{खारे पानी} का इससे छोटा फैलाव है, जो अक्सर किसी
\omterm{def:g2:water-on-earth:ocean}{महासागर} के किनारे होता है।
\end{solution}

\begin{solution}{exo:g2:water-on-earth:5}
नहीं: \omterm{def:g2:water-on-earth:ocean}{महासागर} का पानी \omterm{def:g2:water-on-earth:salt-fresh}{खारा पानी} है, और उसे पीने से शरीर से पानी निकल जाता
है। नाविक को \omterm{def:g2:water-on-earth:salt-fresh}{मीठा पानी} साथ ले जाना पड़ता है।
\end{solution}

\begin{solution}{exo:g2:water-on-earth:6}
उदाहरण के लिए: \omterm{def:g1:landscapes-around-us:water}{नदियों} में, \omterm{def:g1:landscapes-around-us:water}{झीलों} में, \omterm{def:g1:landscapes-around-us:landform}{पर्वतों} और अंटार्कटिका की बर्फ़ और
हिम में, ज़मीन के नीचे, बारिश में।
\end{solution}

\begin{solution}{exo:g2:water-on-earth:7}
\omterm{def:g2:water-on-earth:salt-fresh}{खारा पानी}: 96 बाल्टियाँ। \omterm{def:g2:water-on-earth:salt-fresh}{मीठा पानी}: $2 + 2 = 4$ बाल्टियाँ।
(ठीक-ठीक कहें तो ज़मीन के नीचे का कुछ पानी भी खारा होता है; चित्र उसे \omterm{def:g2:water-on-earth:salt-fresh}{मीठे
पानी} के साथ गिनता है।)
\end{solution}

\begin{solution}{exo:g2:water-on-earth:8}
\omterm{def:g2:water-on-earth:ocean}{महासागर} (96 बाल्टियाँ), फिर बर्फ़ और हिम (2 बाल्टियाँ), फिर \omterm{def:g1:landscapes-around-us:water}{नदियाँ} (एक
बाल्टी से बहुत कम)।
\end{solution}

\begin{solution}{exo:g2:water-on-earth:9}
वह ज़मीन में रिस सकता है; वह ज़मीन पर बहकर किसी नाले या \omterm{def:g1:landscapes-around-us:water}{नदी} में जा सकता
है; धूप निकलने पर वाष्प बनकर वापस हवा में जा सकता है।
\end{solution}

\begin{solution}{exo:g2:water-on-earth:10}
तसले की तली में नमक बचता है। पानी का वाष्पन हो गया है: वह जलवाष्प बनकर
हवा में चला गया है।
\end{solution}

\begin{solution}{exo:g2:water-on-earth:11}
उदाहरण के लिए: बूँद छत से बहकर परनाले और नाली में जाती है, फिर एक नाले और
एक \omterm{def:g1:landscapes-around-us:water}{नदी} में; \omterm{def:g1:landscapes-around-us:water}{नदी} उसे \omterm{def:g2:water-on-earth:ocean}{समुद्र} तक ले जाती है; किसी धूप वाले दिन \omterm{def:g2:water-on-earth:ocean}{समुद्र} से
उसका वाष्पन होता है, वह ऊपर हवा में उठती है और एक बादल का हिस्सा बन जाती
है।
\end{solution}
